\subsectionaddtolist{۳.۱ حملات علیه TLS}

\subsubsection*{آ}
در پروتکل TLS ، مهاجم می‌تواند پیام \lr{ClientKeyExchange} را رهگیری و تغییر دهد، اما این تلاش برای ربایش با شکست مواجه می‌شود. دلیل آن این است که تمامی پیام‌های مبادله شده در طول \lr{Handshake} ، در انتها توسط پیام \lr{Finished} احراز اصالت می‌شوند. پیام \lr{Finished} شامل یک MAC است که روی هش کل تاریخچه پیام‌های قبلی محاسبه می‌شود. از آنجا که مهاجم کلید اصلی و مخفی را ندارد، نمی‌تواند پیام \lr{Finished} معتبری از جانب کلاینت تولید کند.
این دستکاری دقیقا در {آخرین مرحله \lr{Handshake}} ، یعنی زمانی که سرور پیام \lr{Finished} کلاینت را دریافت کرده و هش تاریخچه را با هش محاسبه‌شده توسط خودش مطابقت می‌دهد، آشکار شده و ارتباط فورا قطع می‌گردد.

\subsubsection*{ب}
\begin{enumerate}
    \item  مهاجم برای آغاز این حمله باید پیام \lr{ClientHello} را که در همان ابتدای \lr{Handshake} از سمت کلاینت فرستاده می‌شود رهگیری کند. او فیلدها و اکستنشن‌های مربوط به نسخه‌های جدید (مثل \lr{Supported Versions}) را حذف کرده یا تغییر می‌دهد تا سرور تصور کند کلاینت تنها قادر به پشتیبانی از نسخه‌های قدیمی‌تر است.
    \item  در صورت موفقیت اولیه، نسخه پروتکل توافق‌شده به یک نسخه قدیمی‌تر (مثل \lr{TLS 1.2} یا \lr{TLS 1.1}) تنزل می‌یابد. در پی آن، پارامترهای امنیتی و \lr{Cipher Suites} نشست نیز ضعیف‌تر می‌شوند؛ مثلا به جای تبادل کلید امن با خاصیت \lr{Forward Secrecy} ، ممکن است از \lr{RSA} استاتیک و الگوریتم‌های هش ضعیف‌تری استفاده شود.
    \item  مکانیزم هوشمندانه‌ای برای جلوگیری از این حمله تعبیه شده است. اگر سروری از \lr{TLS 1.3} پشتیبانی کند اما به دلیل درخواست کلاینت (که دستکاری شده) مجبور به توافق روی \lr{TLS 1.2} شود، یک رشته بایت ثابت و مشخص (\lr{Magic String}) را در ۸ بایت آخر فیلد \lr{ServerHello.Random} قرار می‌دهد. کلاینت مدرن (که از \lr{TLS 1.3} پشتیبانی می‌کند) پس از دریافت این پیام، آن ۸ بایت را بررسی می‌کند. با دیدن این رشته ثابت، کلاینت متوجه می‌شود که سرور قابلیت نسخه‌های بالاتر را داشته و حمله \lr{Downgrade} در میانه راه رخ داده است، لذا ارتباط را \lr{Abort} می‌کند.
\end{enumerate}

\subsectionaddtolist{۳.۲ پیکربندی IPsec برای ارتباطات سازمانی}

\subsubsection*{آ}
\begin{enumerate}
    \item  استفاده از {ESP} مناسب‌تر است. از آنجا که ترافیک از اینترنت عمومی عبور می‌کند، برای حفاظت از داده‌های سازمانی علاوه بر احراز اصالت، به {محرمانگی} نیز نیاز مبرم داریم تا امکان شنود وجود نداشته باشد. پروتکل AH داده‌ها را رمزگذاری نمی‌کند.
    \item  باید از {\lr{Tunnel Mode}} استفاده شود. در ارتباط بین دو درگاه \lr{VPN} (مانند روترهای لبه شبکه اصفهان و تهران)، بسته‌های اصلی که حاوی IP های نامعتبر داخلی هستند، باید به طور کامل درون یک IP هدر جدید (با آدرس‌های پابلیک درگاه‌ها) کپسوله شوند تا در بستر اینترنت قابل مسیریابی باشند.
    \item 
    ترتیب فیلدها از چپ به راست به این صورت است:

    \lr{[New IP] | [ESP Header] | [Orig IP] | [TCP/UDP] | [Data] | [ESP Trailer] | [ESP Auth]}
    
    {بخش‌های رمزگذاری شده:} از ابتدای \lr{[Orig IP]} تا انتهای \lr{[ESP Trailer]}.

    {بخش‌های احراز اصالت شده:} از ابتدای \lr{[ESP Header]} تا انتهای \lr{[ESP Trailer]}.
\end{enumerate}

\subsubsection*{ب}
\begin{enumerate}
    \item  پروتکل \lr{AH} روی کل بسته از جمله بخش‌هایی از \lr{IP} هدر خارجی (مثل آدرس \lr{IP} مبدأ) هش محاسبه کرده و آن را احراز اصالت می‌کند. روتر \lr{NAT} برای عبور دادن ترافیک به اینترنت، آدرس \lr{IP} مبدأ را به آدرس پابلیک هتل تغییر می‌دهد. این تغییر باعث نامعتبر شدن هش \lr{AH} در سمت گیرنده شده و بسته دور انداخته می‌شود.
    \item  پروتکل \lr{ESP} در لایه شبکه کار می‌کند و برخلاف \lr{TCP/UDP} دارای شماره پورت نیست. روترهای \lr{NAT} (به طور خاص \lr{PAT} ) برای نگاشت و تفکیک ترافیک چندین کاربر داخلی روی یک \lr{IP} پابلیک، به پورت‌ها وابسته‌اند. بدون پورت، روتر هتل نمی‌داند بسته‌های بازگشتی از سرور شرکت را باید به کدام دستگاه در شبکه داخلی هتل تحویل دهد.
    \item  برای رفع این مشکل، از ویژگی \lr{NAT Traversal (NAT-T)} استفاده می‌شود. \lr{NAT-T} کل بسته \lr{ESP} را درون یک هدر استاندارد \lr{UDP} (معمولا پورت ۴۵۰۰) کپسوله می‌کند. با اضافه شدن پورت‌های مبدأ و مقصد \lr{UDP}، روتر \lr{NAT} به راحتی می‌تواند عملیات ترجمه آدرس را روی بسته انجام دهد.
    \item  در این سناریو نیز از {\lr{Tunnel Mode}} استفاده می‌شود. زیرا کلاینت (کارمند) یک آدرس \lr{IP} مجازی متعلق به شبکه داخلی شرکت دریافت می‌کند. بسته‌های تولید شده با این \lr{IP} مجازی باید در \lr{Tunnel Mode} درون یک بسته جدید با \lr{IP} فیزیکی شبکه هتل کپسوله شوند.
\end{enumerate}

\subsubsection*{ج}
\begin{enumerate}
    \item  این حمله {موفقیت‌آمیز نخواهد بود}. پروتکل IPsec به صورت ذاتی برای مقابله با ارسال مجدد بسته‌های قانونی محافظت شده است.
    \item  این پروتکل از یک {\lr{Sequence Number}} همراه با یک {پنجره ضد ارسال مجدد (\lr{Sliding Window})} استفاده می‌کند. فرستنده به هر بسته یک شماره توالی افزایشی و یکتا اختصاص می‌دهد. گیرنده نیز پنجره‌ای از شماره توالی‌های قابل قبول نگه می‌دارد. اگر بسته‌ای دریافت شود که شماره توالی آن تکراری بوده یا از کران پایین پنجره قدیمی‌تر باشد (مانند بسته‌ای که یک ساعت پیش ارسال شده)، بلافاصله آن را به عنوان یک حمله \lr{Replay} شناسایی کرده و دور می‌ریزد.
    \item اگر \lr{Anti-Replay Window} به اشتباه غیرفعال شده یا اندازه‌اش آن‌قدر بزرگ تنظیم شده باشد که توالی‌های یک ساعت پیش هنوز از پنجره خارج نشده باشند، گیرنده بسته تکراری را به عنوان یک بسته معتبر می‌پذیرد. در نتیجه، سیستم بانکی شرکت همان دستور انتقال وجه را مجددا اجرا کرده و تراکنش بانکی {تکرار می‌شود} که خسارت مالی سنگینی به همراه خواهد داشت.
\end{enumerate}